\begin{quote}
\hfill    ``{\em What I cannot create, I do not understand}.''

$~$ \hfill --- Richard Feynman   
\end{quote}


Representation learning: Identify the low-dimensional distribution {\em and} transform it to a desired compact and structured form (e.g. making the distribution more piecewise linear and independent). 

From nonlinear to linear and from dependent to independent, unrolled optimization that maximizes rate reduction or information gain to promote linearity and independence, via iterative compression and transformation. This chapter connects deep networks to the roots of iterative optimization via gradient descent. 

\begin{definition}[(Open-Loop) Representation]
    Given data \(\vX\), an \textbf{open-loop representation} is a function \(f\) such that the \textit{features} \(\vZ = f(\vX)\) are \textit{compact and structured}.
\end{definition}



\section{Informative Representations}
\subsection{The Principle of Maximal Information Gain}
Introducing and formulating the principle of {\em  information gain}. Transform the data distribution to maximize information gain.

More specifically rate reduction in the important case of mixture of low-dimensional Gaussians. Geometric properties of the global optimization landscape of the associated rate reduction objective. 

\section{White-box Deep Networks via Unrolled Optimization}

\subsection{Mystery about black-box deep networks}
Without breaking the flow, raise the mystery about representation learning via deep neural networks. How to understand the role of all the empirically designed deep neural networks, which are largely black boxes. 

\subsection{Deep networks for sparse representation}
Briefly recapping the sparse Dictionary Learning problem in Chapter \ref{ch:linear-independent}, fine-tuning its unrolled optimization leads to the Learned ISTA (LISTA) architecture. The sparseLand-type networks from unrolling convolutional sparse recovery... 

\subsection{Deep networks for rate reduction}
More formally, by unrolling gradient descent for the rate reduction objective, we derive  \href{https://jmlr.org/papers/v23/21-0631.html}{the white-box deep network ReduNet} (and its connections to ResNet, ResNEXT, and CNNs).


\section{White-box Transformers}
By unifying rate reduction, denoising and diffusion, and transformers, we derive  \href{https://arxiv.org/abs/2306.01129}{the white-box transformer CRATE}.

\section{New Deep Architectures}
New variants of the CRATE networks that are more efficient or more effective. 